\documentclass{article}
\usepackage{corl_2026}
\usepackage{amsmath,amssymb,bm,booktabs,array,graphicx}
\title{Supplementary Material: \\
       Receding-Horizon Target Streams for Learned Fixed-Wing Maneuvering}
\author{Anonymous Authors}

\begin{document}
\maketitle

%===============================================================================
\appendix
\section{Policy Architecture and PPO Training Details}
\label{app:architecture}

\subsection{GRU Actor-Critic Architecture}

The flight skill is a recurrent actor-critic policy with a Gated Recurrent Unit
(GRU) hidden state.  The architecture is identical for both the Euler-angle and
quaternion-conditioned variants; only the observation dimensionality differs
($d_{\mathrm{obs}} = 22$ for the Euler variant, $d_{\mathrm{obs}} = 21$ for the
quaternion variant).

\begin{itemize}
  \item \textbf{Observation embedding:} Dense layer, 128 units, orthogonal
        initialization ($\sqrt{2}$), ReLU activation.
  \item \textbf{Recurrent core:} Single-layer GRU with 128 hidden units,
        scanned over time via \texttt{flax.linen.scan}.
  \item \textbf{Post-GRU:} Dense layer, 256 units $\rightarrow$ LayerNorm
        $\rightarrow$ ReLU.
  \item \textbf{Actor heads:} Five independent categorical heads, one per
        action channel.  Each head projects from a shared 128-unit hidden
        layer to the corresponding logit dimension.  All heads use
        orthogonal initialization~($0.01$) except the speed-brake head, which
        uses a zero-mean bias initialization.
  \item \textbf{Critic head:} Dense 128 $\rightarrow$ ReLU $\rightarrow$
        Dense~1 (orthogonal~$1.0$).
\end{itemize}

\subsection{Observation Components}

The 21-dimensional quaternion-conditioned observation vector contains:

\begin{itemize}
  \item \textbf{Attitude error:} vector part of the shortest-rotation
        quaternion error $q_e$ (3 dimensions).
  \item \textbf{Body-frame target direction:} target heading and pitch
        transformed into a body-frame unit vector (3 dimensions).
  \item \textbf{Speed and altitude:} normalized speed error,
        normalized current speed, normalized altitude (3 dimensions).
  \item \textbf{Aircraft state:} body-axis angular rates $P,Q,R$,
        angle-of-attack and sideslip features $\sin\alpha,\cos\alpha,\beta$
        (5~dimensions).
  \item \textbf{Previous action:} one-hot encoding of the last five-channel
        discrete action (5~dimensions).
  \item \textbf{Additional:} remaining dimensions encode task context and
        normalization constants.
\end{itemize}

\subsection{PPO Hyperparameters}

Table~\ref{tab:ppo-hparams} lists the Proximal Policy Optimization
hyperparameters used for training the base skill.  The same configuration
is used for the energy-aware fine-tuning stage unless noted otherwise.

\begin{table}[h]
  \centering
  \caption{PPO hyperparameters.}
  \label{tab:ppo-hparams}
  \begin{tabular}{lr}
    \toprule
    \textbf{Parameter} & \textbf{Value} \\
    \midrule
    Number of parallel environments  & 1000 \\
    Rollout steps per update         & 1000 \\
    Total environment steps          & $1 \times 10^9$ \\
    PPO epochs per update            & 16 \\
    Number of minibatches            & 5 \\
    Discount factor $\gamma$         & 0.99 \\
    GAE parameter $\lambda$          & 0.95 \\
    PPO clipping $\epsilon$          & 0.2 \\
    Value-function clipping $\epsilon_{\!v}$ & 0.2 \\
    Initial learning rate (Adam)     & $3 \times 10^{-4}$ \\
    Initial entropy coefficient      & $1 \times 10^{-3}$ \\
    Entropy coefficient range        & $[5 \times 10^{-4},\; 2 \times 10^{-2}]$ \\
    Optimizer                        & Adam ($\beta_1{=}0.9$, $\beta_2{=}0.999$) \\
    Gradient norm clipping           & 0.5 \\
    \bottomrule
  \end{tabular}
\end{table}

\subsection{Energy-Aware Fine-Tuning Curriculum}

The base skill is trained on a horizontal and mild-attitude tracking
distribution.  To extend capability toward vertical arcs, we fine-tune
with a curriculum that introduces progressively more aggressive vertical
tasks while replaying original tasks to prevent catastrophic forgetting.

The curriculum defines two task regimes:
\begin{enumerate}
  \item \textbf{Horizontal retention tasks} (50\% probability): heading
        steps ($\pm 20^\circ$, $\pm 45^\circ$), pitch steps ($\pm 10^\circ$),
        level flight, level circles ($R = 3000,5000$~m), S-curves, figure-eights,
        altitude steps ($\pm 1000$~m).
  \item \textbf{Vertical curriculum tasks} (50\% probability): pull-ups
        ($15^\circ$--$30^\circ$, $R = 2000$--$8000$~m), loop-plane arcs
        ($60^\circ$--$150^\circ$, $R = 8000$--$15000$~m).
\end{enumerate}
The vertical stage advances when the aircraft successfully completes a
configurable number of consecutive tracking attempts, progressively
exposing the policy to larger arc angles.

The residual specialist for large loop-plane arcs ($> 80^\circ$) is
trained separately from the base skill.  It receives an augmented
observation that includes the current loop-plane phase angle and a gate
signal that activates progressively between $80^\circ$ and $180^\circ$
along the arc.  The residual policy outputs additive logit deltas to the
base policy's action distribution, clipped to
$[- \text{clip}, +\text{clip}]$ where $\text{clip} = 1.25$.  The
residual network uses a smaller GRU (64 hidden units, FC~dim~96) and is
trained with the same PPO configuration as the base policy.

%===============================================================================
\section{Reward Terms}
\label{app:rewards}

The energy-aware fine-tuning reward combines ten terms:

\begin{align}
  r &= 0.68\, r_{\mathrm{att}} \;+\; 0.24\, r_{\mathrm{speed}}
       \;+\; r_{\mathrm{energy}} \;+\; r_{\mathrm{climb}}
       \;+\; r_{\mathrm{alt}} \;+\; r_{\alpha\beta}
       \;+\; r_{G} \;+\; r_{\mathrm{smooth}} \;+\; r_{\mathrm{replay}}
       \;+\; r_{\mathrm{alive}} ,
\end{align}
clipped to $[-2.0, 1.25]$.  The terms are as follows.

\paragraph{Attitude tracking ($r_{\mathrm{att}}$).}
\begin{equation}
  r_{\mathrm{att}} = \exp\!\bigl(-(\theta / \theta_0)^2\bigr),
\end{equation}
where $\theta$ is the geodesic angle between current and target attitude
quaternions and $\theta_0$ is a scaling constant.

\paragraph{Speed tracking ($r_{\mathrm{speed}}$).}
\begin{equation}
  r_{\mathrm{speed}} = \exp\!\bigl(-(\Delta v_{\mathrm{under}} / s_0)^2
                         - (\Delta v_{\mathrm{over}} / s_1)^2\bigr),
\end{equation}
with separate penalties for under-speed and over-speed relative to the
target airspeed.

\paragraph{Energy retention ($r_{\mathrm{energy}}$).}
\begin{equation}
  r_{\mathrm{energy}} = -\bigl(0.06\, \Delta E_{\mathrm{step}}^2
                        + 0.04\, \Delta E_{\mathrm{task}}^2\bigr)
                        \cdot g_{\mathrm{vertical}},
\end{equation}
where $\Delta E_{\mathrm{step}}$ is the per-step specific-energy change,
$\Delta E_{\mathrm{task}}$ is the change since task start, and
$g_{\mathrm{vertical}}$ is a gate that activates on vertical-curriculum
tasks.

\paragraph{Climb progress ($r_{\mathrm{climb}}$).}
\begin{equation}
  r_{\mathrm{climb}} = 0.20 \cdot \tanh(\Delta h / 350)
                       \cdot [v_t > v_{\min}]
                       \cdot g_{\mathrm{climb}},
\end{equation}
rewarding altitude gain when the aircraft maintains adequate airspeed.

\paragraph{Altitude hold ($r_{\mathrm{alt}}$).}
A penalty proportional to the squared deviation from a reference
altitude band, active during non-vertical tasks.

\paragraph{$\alpha$/$\beta$/\textit{G} penalties ($r_{\alpha\beta}$, $r_G$).}
\begin{align}
  r_{\alpha\beta} &= -\,0.08\,(\max(0, |\alpha| - \alpha_{\max}))^2
                     - 0.04\,(\max(0, |\beta| - \beta_{\max}))^2 ,\\
  r_{G}           &= -\,0.12\,(\max(0, G - G_{\max}))^2 .
\end{align}
These terms penalize excursions beyond safe aerodynamic and structural
envelopes ($\alpha_{\max} \approx 25^\circ$, $\beta_{\max} \approx 5^\circ$,
$G_{\max} \approx 8.5$).

\paragraph{Action smoothness ($r_{\mathrm{smooth}}$).}
\begin{equation}
  r_{\mathrm{smooth}} = -0.025 \cdot
    (0.5\,\Delta\mathrm{thr} + \Delta\mathrm{ele}
     + 0.35\,\Delta\mathrm{ail} + 0.35\,\Delta\mathrm{rud}
     + 0.5\,\Delta\mathrm{sb}),
\end{equation}
penalizing large actuator changes between consecutive control steps.

\paragraph{Alive bonus ($r_{\mathrm{alive}}$).}
A constant $+0.02$ per step awarded for survival.

\paragraph{Replay term ($r_{\mathrm{replay}}$).}
Additional geometry-alignment terms (velocity-tangent, nose-tangent,
nose-velocity, wing-plane errors) active during loop-like phases to
enforce proper maneuver-plane orientation.

%===============================================================================
\section{Target-Stream Generator and Search Space}
\label{app:target-stream}

\subsection{Target-Stream Definition}

Given a reference maneuver $\mathcal{M}$ (expressed as a sequence of
3D waypoints) and the current aircraft state $x_t$, a target-stream
generator $G$ produces a local target sequence
\begin{equation}
  \tau_{t:t+H} = \bigl\{(\psi_k^\star, \theta_k^\star, \phi_k^\star, v_{t,k}^\star)\bigr\}_{k=t}^{t+H},
\end{equation}
where each $\tau_k$ is the desired heading, pitch, roll, and airspeed
that the flight skill is conditioned on.

\subsection{Lookahead-Based Stream Generation}

The generator $G$ is parameterized by a scalar lookahead distance $L$
and a target airspeed $v_t^\star$.  At each step, the generator
projects the aircraft position onto the reference path, advances along
the arc by $L$ meters, and computes the heading and pitch from the
relative position error to the lookahead point:
\begin{align}
  \psi_k^\star &= \arctan2(e_{\mathrm{la}} - e_{\mathrm{ac}},
                            n_{\mathrm{la}} - n_{\mathrm{ac}}), \\
  \theta_k^\star &= \arctan2(a_{\mathrm{la}} - a_{\mathrm{ac}},
                              \sqrt{(n_{\mathrm{la}} - n_{\mathrm{ac}})^2
                                  + (e_{\mathrm{la}} - e_{\mathrm{ac}})^2}),
\end{align}
where $(n_{\mathrm{la}}, e_{\mathrm{la}}, a_{\mathrm{la}})$ is the
lookahead point in world coordinates.  Roll target $\phi_k^\star$ is
set to zero (coordinated flight).

\subsection{Offline Stream Selection}

The optimal target-stream parameters for each maneuver class are
determined through \textbf{offline closed-loop evaluation}.  For each
candidate pair $(L, v_t^\star)$ in the lattice, a full closed-loop
rollout of the frozen policy is executed on the reference trajectory.
The candidate that minimizes the cross-track error (CTE-P90) over the
entire trajectory is selected as the task-optimal parameterization.

This offline sweep reveals that the optimal parameters are
\textbf{task-dependent}: horizontal and mild 3D maneuvers
(S-curve, figure-eight, helix) prefer shorter lookahead and lower
target speed $(600, 220)$, while the $90^\circ$ vertical pull-up
requires longer lookahead and higher target speed $(1500, 280)$.
The default stream $(1000, 250)$ used as a comparison baseline is
optimal for none of the tested tasks, empirically motivating the
need for task-adaptive stream selection.

\subsection{Candidate Lattice}

The offline sweep evaluates three candidate parameter pairs:
\begin{equation}
  \mathcal{L} = \{(600, 220),\; (1000, 250),\; (1500, 280)\},
\end{equation}
where each pair is $(L\ \text{[m]},\ v_t^\star\ \text{[m/s]})$.
The default stream used as a baseline is $(L, v_t^\star) = (1000, 250)$.

%===============================================================================
\section{Classical Autopilot Baseline Details}
\label{app:classical}

\subsection{Three-Layer Architecture}

The classical autopilot baseline consists of three layers inspired by
the JSBSim F-16 flight control system:

\begin{enumerate}
  \item \textbf{Inner SAS (Stability Augmentation System).}
        Pitch-axis C* blending (pitch-rate + normal-acceleration feedback),
        roll-axis rate damping, yaw-axis damper with sideslip suppression.
        Alpha limiting at $|\alpha| > 25^\circ$.

  \item \textbf{Middle PID Attitude Controller.}
        Converts attitude errors (pitch, roll) into desired angular-rate
        commands, and speed error into a throttle command.

  \item \textbf{Outer NLGL (Nonlinear Guidance Law).}
        Projects the aircraft position onto the reference path, looks ahead
        by a distance $L_1$, and computes desired heading, pitch, and
        airspeed from the relative position error.
\end{enumerate}

\subsection{PID Gains and Tuning}

Table~\ref{tab:pid-gains} reports the empirically tuned PID gains used
in all classical-baseline experiments.  Gains were selected through a
systematic sweep to achieve stable straight-and-level flight (verified
over 400+ control steps / 80~seconds).

\begin{table}[h]
  \centering
  \caption{PID gains for the classical autopilot.  Sign convention:
           negative elevator $\rightarrow$ nose-up.  Elevator trim
           $= -0.025$ for level flight at 5000~m / 250~m/s.}
  \label{tab:pid-gains}
  \begin{tabular}{lcccc}
    \toprule
    \textbf{Channel} & \textbf{$k_p$} & \textbf{$k_i$} & \textbf{$k_d$} & \textbf{Integrator Max} \\
    \midrule
    Pitch $\rightarrow$ elevator  & 0.3 & 0.03 & 0.15 & 0.4 \\
    Roll $\rightarrow$ aileron    & 0.4 & 0.02 & 0.10 & 0.3 \\
    Speed $\rightarrow$ throttle  & 0.015 & 0.005 & 0.0 & 0.3 \\
    Yaw $\rightarrow$ rudder      & 0.3 & 0.0 & 0.05 & 0.1 \\
    \bottomrule
  \end{tabular}
  \vspace{4pt}
  \footnotesize
  A gravity-compensation feedforward term $-k_{\!f\!f} \sin(\theta_{\mathrm{target}})$
  with $k_{\!f\!f} = 0.03$ is added to the elevator command to provide
  sustained nose-up bias at high pitch angles.
\end{table}

\subsection{Control Rates and Comparison Rationale}

Both the classical autopilot and the RL policy operate at the identical
control frequency of \textbf{5~Hz} (one action per 10 simulation steps at
50~Hz physics integration).  The classical inner-loop PID and outer-loop
NLGL are updated synchronously at this rate.  We also tested a
\textbf{multi-rate variant} with the PID inner loop at 50~Hz and NLGL
outer loop at 5~Hz.  The multi-rate variant showed \emph{worse}
survival (43 steps vs.\ 119 steps on the $90^\circ$ pull-up), because
the PID integral accumulates within the NLGL update interval without
fresh guidance targets, leading to windup.  We therefore report the
synchronous 5~Hz variant as the representative classical baseline.

\subsection{Failure Analysis on the 90-Degree Pull-Up}

The classical autopilot fails on the $90^\circ$ vertical pull-up
($R = 10000$~m) after 119 control steps (23.8~s), having reached an
altitude of 6985~m before triggering the environment's overload
protection.  The failure proceeds through three phases:

\begin{enumerate}
  \item \textbf{Integral saturation (0--14~s).}
        The slow actuator dynamics ($e\ell \leftarrow 0.9\,e\ell +
        0.1\,\mathrm{cmd}\cdot 45^\circ$, $\tau \approx 1.8$~s) cause
        the actual pitch to overshoot the NLGL target.  The PID integral
        saturates at $-0.4$, commanding nose-down while the gravity
        feedforward simultaneously commands nose-up.
  \item \textbf{PIO divergence (22.0--23.6~s).}
        The conflicting commands induce a pilot-induced oscillation:
        pitch oscillates between $38^\circ$ and $53^\circ$, angle of
        attack surges from $10^\circ$ to $18.4^\circ$, and G-load
        reaches $9.6$~g.
  \item \textbf{Overload termination (23.8~s).}
        The environment's overload protection ($G \geq 10$~g, or
        $|\alpha| \geq 25^\circ$) terminates the episode.
\end{enumerate}

The RL optimal variant completes the same maneuver in 336 steps (67.2~s),
reaching 14736~m with a peak G-load of 7.1~g and peak angle of attack
of $3.3^\circ$.  The fundamental difference is the RL policy's implicit
feedforward compensation (learned through the GRU hidden state), which
anticipates the required elevator trim before large errors develop,
whereas the classical PID is purely reactive and cannot overcome the
combined effects of actuator lag and the airframe's relaxed static
stability at this flight condition.

\end{document}
